Esta práctica buscó determinar la densidad 
en diversos cuerpos sólidos: un corcho, 
dos cilindros metálicos, un cilindro hueco y una 
moneda. Se utilizaron equipos de medición de precisión, 
tales como una balanza granataria, un vernier y un micrómetro 
para obtener las masas y dimensiones de los objetos. Los datos 
recolectados se procesaron estadísticamente para obtener los valores 
necesarios. Posteriormente, se realizó el cálculo de propagación de 
incertidumbres mediante el método de derivadas parciales para validar 
la confiabilidad de las mediciones. Se concluye que la densidad de los 
objetos varía significativamente; el corcho presenta el valor 
menor, mientras que los cilindros metálicos muestran los valores 
más altos, resultados 
que son consistentes con las propiedades intensivas de cada material. 